% ============================================================
% Come si usano, argomento 19: decisioni bayesiane e classificazione binaria
% ============================================================
\argomento{Decisioni bayesiane e classificazione binaria}{Formule a p.~\pageref*{f:regoladecisione}}[Decisioni
  bayesiane]

\begin{scelta}{Come scegliere la formula}
\renewcommand{\arraystretch}{1.25}
\begin{tabular}{@{}P{0.55\linewidth}P{0.39\linewidth}@{}}
  \toprule
  \textbf{Se il testo\ldots} & \textbf{si usa\ldots} \\
  \midrule
  chiede la regola che sbaglia meno, con le probabilità a priori note & la regola MAP \vv{regolamap} \\
  dice che le ipotesi sono equiprobabili, o non dà le probabilità a priori & la regola ML \vv{regolaml} \\
  dà costi diversi per errori diversi & il rischio bayesiano \vv{rischio} \\
  chiede quanto spesso la regola sbaglia & \vv{perrore} \\
  ha dati binari, da due sorgenti con diverse probabilità di 1 & \vv{testhamming} \\
  ha misure gaussiane con due medie possibili & \vv{testmedia} \\
  \bottomrule
\end{tabular}

Il procedimento è sempre lo stesso: si scrivono le verosimiglianze dei dati sotto le due ipotesi, le si confronta
pesate con le probabilità a priori e si semplifica fino a una condizione su una statistica sufficiente, come il numero
di uni o la media dei dati.
\end{scelta}

\scheda{regoladecisione}{
  \usoquando{Si deve scegliere fra $M$ spiegazioni dei dati: quale simbolo è stato trasmesso, se un messaggio è posta
    indesiderata, l'esito di un test medico.}
  \usocome{Si dice, per ogni dato possibile, quale ipotesi si sceglie: le regioni $\Omega_i$ non si sovrappongono e
    coprono tutti i dati.}
  \usoesempio{Un bit $\pm 1$ ricevuto con rumore, $x \in \R$: $D(x) = 1$ se $x > 0$ e $D(x) = 2$ altrimenti, cioè
    $\Omega_1 = (0, +\infty)$ e $\Omega_2 = (-\infty, 0]$.}
}

\scheda{rischio}{
  \usoquando{Gli errori non costano tutti uguale: scambiare un messaggio importante per posta indesiderata è peggio del
    contrario.}
  \usocome{Ogni costo si pesa con la probabilità di decidere $i$ quando è vera $H_j$. Con costi 0 per le decisioni
    giuste e 1 per quelle sbagliate il rischio è la probabilità di errore.}
  \usoesempio{Test medico, $H_1$ malato e $H_2$ sano: un malato non visto costa $C_{21} = 10$, un falso allarme
    $C_{12} = 1$. Con $P\{D = 2, H_1\} = 0{,}03$ e $P\{D = 1, H_2\} = 0{,}05$: $R = 10 \cdot 0{,}03 + 1 \cdot 0{,}05
    = 0{,}35$, mentre $P(e) = 0{,}08$.}
}

\scheda{perrore}{
  \usoquando{Si valutano le prestazioni di una regola: quanto spesso sbaglia, in tutto e sotto ciascuna ipotesi.}
  \usocome{Si calcola la probabilità di cadere nella regione sbagliata sotto ciascuna ipotesi, poi si pesano le due
    con le probabilità a priori (probabilità totale, \vv{totale}).}
  \usoesempio{Test sul peso di Hamming con dieci bit: $P(e \mid H_1) \simeq 0{,}055$ e $P(e \mid H_2) \simeq 0{,}121$,
    quindi $P(e) \simeq \frac{1}{2}(0{,}055 + 0{,}121) \simeq 0{,}088$.}
}

\scheda{regolamap}{
  \usoquando{Si conoscono le probabilità a priori delle ipotesi e si vuole sbagliare il meno possibile.}
  \usocome{Si calcolano le due verosimiglianze nei dati osservati, se ne fa il rapporto e lo si confronta con $\eta =
    \frac{P(H_2)}{P(H_1)}$. Spesso conviene prendere i logaritmi dei due membri: il verso del confronto non cambia,
    perché il logaritmo è crescente.}
  \usoattenzione{Le probabilità a priori spostano la soglia: per scegliere un'ipotesi rara servono dati che la
    favoriscano nettamente.}
  \usoesempio{Sorgenti con $p_1 = 0{,}6$ e $p_2 = 0{,}2$, dieci bit, $P(H_1) = 0{,}2$: $\eta = 4$, e $L = 3^w
    \cdot 0{,}5^{10 - w} > 4$ solo con $w \geq 5$ uni (con ipotesi equiprobabili ne bastavano 4). $P(e) \simeq
    0{,}059$, contro $0{,}108$ tenendo la soglia 4.}
}

\scheda{regolaml}{
  \usoquando{Le ipotesi sono equiprobabili, o le probabilità a priori non si conoscono: si sceglie l'ipotesi che rende
    più probabili i dati osservati.}
  \usoesempio{Dieci bit con 4 uni, sorgenti con $p_1 = 0{,}6$ e $p_2 = 0{,}2$: $L = 3^4 \cdot 0{,}5^6 = \frac{81}{64}
    \simeq 1{,}27 > 1$, si decide $H_1$; con 3 uni $L \simeq 0{,}21$ e si decide $H_2$.}
}

\scheda{testhamming}{
  \usoquando{I dati sono bit i.i.d. e le due ipotesi differiscono per la probabilità che un bit valga 1.}
  \usocome{Si contano gli uni e si confrontano con la soglia; le probabilità di errore sono code di binomiali
    (\vv{distbinomiale}).}
  \usoesempio{$n = 10$, $p_1 = 0{,}6$, $p_2 = 0{,}2$: $\eta_1 = 10\, \frac{\ln 2}{\ln 6} \simeq 3{,}87$, si decide
    $H_1$ con almeno 4 uni e $P(e) \simeq 0{,}088$. Con soglia 3 o 5 si otterrebbe $0{,}167$ o $0{,}100$.}
}

\scheda{testmedia}{
  \usoquando{Si hanno $n$ misure di una grandezza che può valere $\mu_1$ o $\mu_2$, disturbate da rumore gaussiano.}
  \usocome{Si fa la media delle misure e si guarda se è più vicina a $\mu_1$ o a $\mu_2$. Più misure moltiplicano per
    $\sqrt{n}$ l'argomento di $Q$.}
  \usoesempio{Tensione di $\pm 1$ V con rumore $\sigma = 2$ V: con una misura $P(e) = Q(0{,}5) \simeq 0{,}31$; con 16
    misure $P(e) = Q(2) \simeq 0{,}023$.}
}
